%==================================================================
% Ini adalah bab 5
%==================================================================

\chapter[KESIMPULAN DAN SARAN]{\\ KESIMPULAN DAN SARAN}

\section{Kesimpulan}
Berdasarkan hasil perancangan, implementasi, dan pengujian sistem klasifikasi kelayakan air pendingin (\textit{cooling water}) pada mesin generator W{\"a}rtsil{\"a} 20V34SG menggunakan algoritma \textit{Random Forest}, dapat ditarik beberapa kesimpulan utama sebagai berikut:
\begin{packed_enum}
    \item Sistem klasifikasi berbasis web telah berhasil dibangun dan diimplementasikan menggunakan \textit{framework} Flask untuk backend, HTML/CSS/JS untuk frontend, dan MySQL sebagai pengelola basis data. Sistem ini memudahkan Operator dalam memantau kelayakan air pendingin secara terintegrasi serta mendigitalisasi pencatatan data riwayat pengujian laboratorium.
    \item Penerapan algoritma \textit{Random Forest} terbukti sangat efektif untuk mengklasifikasikan status kelayakan air pendingin menjadi kategori ``Layak'' atau ``Tidak Layak''. Dari pengujian model final pada data uji (\textit{testing set}), diperoleh tingkat akurasi mencapai \textbf{98,48\%}, presisi sebesar \textbf{100,00\%}, recall sebesar \textbf{97,56\%}, dan F1-score sebesar \textbf{98,77\%}.
    \item Hasil pengujian menunjukkan nilai \textit{False Positive} sebesar \textbf{0,00\%}, yang berarti sistem tidak pernah salah mengklasifikasikan air pendingin yang sebenarnya ``Tidak Layak'' diprediksi sebagai ``Layak''. Hal ini sangat penting untuk mencegah penggunaan air pendingin yang tidak memenuhi syarat mutu standar pada mesin generator W{\"a}rtsil{\"a}, sehingga dapat meminimalkan risiko pembentukan kerak (\textit{scaling}) dan korosi pada pipa pendingin mesin pembangkit.
    \item Analisis \textit{Feature Importance} menunjukkan bahwa parameter \textbf{Nitrite} ($NO_2$) merupakan fitur yang paling berpengaruh dalam penentuan keputusan kelayakan air dengan tingkat kepentingan sebesar \textbf{35,0\%}, diikuti oleh \textbf{Specific Conductivity} (\textbf{20,0\%}), dan \textbf{pH} (\textbf{18,0\%}). Hal ini sejalan dengan teori teknik kimia di mana nitrit bertindak sebagai korosi inhibitor utama dalam sistem air pendingin tertutup.
    \item Hasil pengujian fungsionalitas menggunakan metode \textit{Black Box} menunjukkan bahwa seluruh fungsi utama sistem (seperti autentikasi login, input parameter, proses prediksi, penyaringan laporan, ekspor data excel, dan kelola user) berjalan dengan sukses dan sesuai harapan tanpa adanya kegagalan program. Pengujian \textit{White Box} juga mengonfirmasi keandalan alur logika program dalam menangani pengecualian (\textit{exception handling}) saat terjadi kendala pada file model atau koneksi basis data.
\end{packed_enum}

\section{Saran}
Beberapa saran yang dapat diajukan untuk pengembangan dan penyempurnaan sistem ini di masa mendatang adalah sebagai berikut:
\begin{packed_item}
    \item \textbf{Penambahan Jumlah Dataset}: Disarankan untuk terus menambah dan memperbarui entri data hasil laboratorium air pendingin secara berkala agar variasi data semakin kaya, sehingga dapat terus melatih kembali (\textit{retraining}) model \textit{Random Forest} guna meningkatkan akurasi generalisasi dan nilai \textit{recall} hingga mencapai nilai mutlak 100\%.
    \item \textbf{Integrasi Sensor IoT Terdistribusi}: Pengembangan sistem dapat diarahkan pada integrasi dengan sensor \textit{Internet of Things} (IoT) secara \textit{real-time} (seperti sensor pH, temperatur, dan konduktivitas air) untuk mengotomatisasi pengambilan sampel parameter secara langsung dari saluran pipa air pendingin tanpa perlu melakukan input data manual.
    \item \textbf{Sistem Peringatan Dini (Alert System)}: Menambahkan fitur notifikasi otomatis melalui email atau SMS/WhatsApp API yang langsung dikirimkan ke Admin atau Kepala Divisi Pemeliharaan apabila hasil klasifikasi bernilai ``Tidak Layak'', sehingga tindakan korektif (seperti \textit{blowdown} atau penambahan zat inhibitor) dapat segera dilaksanakan sebelum timbul kerusakan mesin.
    \item \textbf{Eksperimen dengan Metode Ensemble Lain}: Melakukan penelitian komparatif menggunakan algoritma \textit{ensemble} lainnya seperti XGBoost, CatBoost, atau optimasi \textit{hyperparameter tuning} menggunakan metode \textit{Genetic Algorithm} (GA) untuk memperbandingkan performa efisiensi komputasi dan akurasi model klasifikasi.
\end{packed_item}